\section*{第 \ref{ch:b1:carbohydrates} 章——糖类}

\begin{solution}{exo:b1:carbohydrates:1}
葡萄糖：己糖，\omterm{def:b1:carbohydrates:monosaccharide}{醛糖}。果糖：己糖，\omterm{def:b1:carbohydrates:monosaccharide}{酮糖}。核糖：戊糖，
\omterm{def:b1:carbohydrates:monosaccharide}{醛糖}。甘油醛：丙糖，\omterm{def:b1:carbohydrates:monosaccharide}{醛糖}。
\end{solution}

\begin{solution}{exo:b1:carbohydrates:2}
成环后的羰基碳连着环氧和一个羟基，该羟基
可重新打开，恢复为醛或酮。在蔗糖中，两个异头碳
都被锁定在\omterm{def:b1:carbohydrates:glycosidic}{糖苷键}中，均无法打开，因此不生成醛，
铜离子也不被还原。
\end{solution}

\begin{solution}{exo:b1:carbohydrates:3}
\omterm{def:b1:carbohydrates:polysaccharide}{淀粉}：葡萄糖，$\alpha(1{\to}4)$ 连接，并带有 $\alpha(1{\to}6)$ 分支
（支链\omterm{def:b1:carbohydrates:polysaccharide}{淀粉}）。\omterm{def:b1:carbohydrates:polysaccharide}{糖原}：结构相同，但分支更多。\omterm{def:b1:carbohydrates:polysaccharide}{纤维素}：葡萄糖，
$\beta(1{\to}4)$ 连接。\omterm{def:b1:carbohydrates:polysaccharide}{几丁质}：N-乙酰葡糖胺，$\beta(1{\to}4)$ 连接。
\end{solution}

\begin{solution}{exo:b1:carbohydrates:4}
C1 向下，C2 向下，C3 向上，C4 向下。$\alpha$-D-半乳糖与之相同，
只是 C4 上的羟基向上。
\end{solution}

\begin{solution}{exo:b1:carbohydrates:5}
在 $\alpha$ 键中，异头羟基朝下，每个残基都能以与前一个相同的角度
接上：链不断转折并闭合成螺旋；在 $\beta$ 键中，
异头羟基朝上，
残基只能翻转后再接上：链呈直线伸展。
碘分子可嵌入直链\omterm{def:b1:carbohydrates:polysaccharide}{淀粉}的螺旋内部，其电子状态随之改变（呈蓝色）；
扁平的带状结构没有空腔。
\end{solution}

\begin{solution}{exo:b1:carbohydrates:6}
碘：只有\omterm{def:b1:carbohydrates:polysaccharide}{淀粉}变为蓝黑色。对其余两种糖用本尼迪特试剂检测：只有
葡萄糖生成红色沉淀，蔗糖仍呈蓝色。为确认蔗糖，先用稀酸加热，
中和后再次用本尼迪特试剂检测，
此时呈阳性。
\end{solution}

\begin{solution}{exo:b1:carbohydrates:7}
每条链含 12 个残基：约 \num{4000} 条链，其中一半位于最外层，
即约 \num{2000} 个末端。直链\omterm{def:b1:carbohydrates:polysaccharide}{淀粉}只有一个末端。要释放 \num{5000} 个葡萄糖，
每个末端每秒释放一个：\omterm{def:b1:carbohydrates:polysaccharide}{糖原}需 \qty{2.5}{s}；直链\omterm{def:b1:carbohydrates:polysaccharide}{淀粉}需 \qty{5000}{s}，
接近一个半小时。
\end{solution}

\begin{solution}{exo:b1:carbohydrates:8}
$\sigma = 0.8\times 25/(2\times 0.5) = \qty{20}{MPa}$：约为微纤丝强度的
五十分之一；\omterm{def:b1:eukaryotic-cell:plantcell}{细胞壁}的安全余量很大，这是必要的，因为微纤丝
仅占\omterm{def:b1:eukaryotic-cell:plantcell}{细胞壁}的四分之一。
\end{solution}

\begin{solution}{exo:b1:carbohydrates:9}
乳糖酶断开半乳糖与葡萄糖之间的 $\beta(1{\to}4)$ 键；
无此酶则乳糖不能被吸收（只有\omterm{def:b1:carbohydrates:monosaccharide}{单糖}能穿过
\omterm{def:b1:body-plans-tissues:epithelium}{上皮}）。在结肠中，细菌将其发酵为酸和气体
（引起痉挛、腹胀），未被消化的糖和酸将水
通过\omterm{def:b1:membranes-transport:osmosis}{渗透作用}吸入肠腔：导致腹泻。
\end{solution}

\begin{solution}{exo:b1:carbohydrates:10}
瘤胃中的细菌和原生生物分泌\omterm{def:b1:carbohydrates:polysaccharide}{纤维素}酶，
并将葡萄糖发酵为短链\omterm{def:b1:lipids:fattyacid}{脂肪酸}；奶牛吸收并氧化这些\omterm{def:b1:lipids:fattyacid}{脂肪酸}，
微生物本身随后被消化。发酵必须先于小肠进行，
使其产物能到达吸收表面，
微生物蛋白也能进入胃和肠道。\omterm{def:b1:carbohydrates:polysaccharide}{纤维素}：
\qty{4}{kg}；发酵掉 \qty{2.4}{kg}；每天 \qty{29}{MJ}，
占奶牛能量的大部分。
\end{solution}

\begin{solution}{exo:b1:carbohydrates:11}
O 型血\omterm{def:b1:cell-unit-of-life:cell}{细胞}不带任何抗原，也不会在任何人体内引发抗体；
但 O 型\omterm{def:b1:mammal-organization:compartments}{血浆}中含有抗 A 和抗 B 的抗体，因此 O 型人
若接受 A 型或 B 型血\omterm{def:b1:cell-unit-of-life:cell}{细胞}，就会发生攻击。A 型人的免疫系统
从未接触过 B 型抗原，会产生抗 B 抗体（针对携带相似糖链的
肠道细菌），但对 A 型抗原有耐受：A 型\omterm{def:b1:mammal-organization:compartments}{血浆}只含抗 B 抗体。
\end{solution}

\begin{solution}{exo:b1:carbohydrates:12}
正确：葡萄糖是几乎所有\omterm{def:b1:cell-unit-of-life:cell}{细胞}的燃料，\omterm{def:b1:carbohydrates:polysaccharide}{纤维素}是地球上
含量最丰富的有机分子，两者都是 D-葡萄糖的聚合物，或即为 D-葡萄糖。
一个键的差别改变了一切：$\alpha$ 键形成的线圈易被酶打开、
可供\omterm{def:b1:cell-unit-of-life:cell}{细胞}燃烧；$\beta$ 键形成的纤维则抵抗酶、水和张力，
而大多数动物根本无法利用。
生物学不是分子的清单，而是其连接方式的几何学
——同一种化学物质，依据单一的立体化学选择，可以是食物，也可以是支架，
而生物的分野在于它们拥有哪些
能够拆解它的酶。
\end{solution}

\begin{solution}{pb:b1:carbohydrates:1}
\textbf{1.} $2^{t-1}$：1、2、4，以及 $2^{11} = 2048$。
\textbf{2.} $2^{12} - 1 = 4095$。
\textbf{3.} $4095\times 13 = \num{53235}$，约 \num{53000}。
\textbf{4.} $\num{53235}\times 162 = \qty{8.6e6}{Da}$；
$\qty{1.43e-17}{g}$。
\textbf{5.} $2048/4095$：约一半。\num{2048} 个非还原末端。
\textbf{6.} 一个。
\textbf{7.} $12\times 1.9 = \qty{23}{nm}$ 为半径，\qty{46}{nm}
直径：约为核糖体的两倍。
\textbf{8.} $100/\num{1.43e-17} = \num{7.0e18}$ 个粒子。
\textbf{9.} $100/162 = \qty{0.62}{mol}$ 残基，即 \qty{111}{g}
游离葡萄糖（水解时加入的水造成了两者之差）。
\textbf{10.} \qty{400}{g}，超过肝脏质量的四分之一。
\textbf{11.} \qty{10}{g/h} $= \qty{2.78}{mg/s} = \qty{1.54e-5}{mol/s} =
\num{9.3e18}$ 个分子每秒。
\textbf{12.} 约 \qty{11}{h}（111 g，按 10 g/h 计）——相当于一夜。主要
供给大脑，大脑耗用 \qty{5}{g/h}。
\textbf{13.} $2048\times\num{7e18} = \num{1.4e22}$ 个末端。
\textbf{14.} $\num{1.4e22}\times 20 = \num{2.9e23}$ 每秒 $=
\qty{0.48}{mol/s} = \qty{86}{g/s}$：约为实际
速率的三万倍。末端从来不是限制因素，酶才是。
\textbf{15.} $\num{9.3e18}/20 = \num{4.6e17}$ 个分子——约
\qty{75}{mg} 酶，仅占肝脏蛋白质的一小部分。
\textbf{16.} 链长为 \num{53000} 个残基：末端数仍为 $\num{7e18}$（每个粒子一个），
少了 \num{2048} 倍；最大速率为 \qty{42}{mg/s}，仍是
实际速率的十五倍。（直链\omterm{def:b1:carbohydrates:polysaccharide}{淀粉}链更长时，余量
会缩小；无分支储备物的真正限制是其不溶性和结晶性，
它们使末端无法被酶接近。）
\textbf{17.} 血糖下降时，激素在数秒内激活磷酸化酶；
酶在每个粒子表面即可找到数千个已暴露的末端，
因此释放速率可一次性提高百倍，
无需等待链解开或
粒子溶解。
\textbf{18.} $0.55/1.0 = \qty{0.55}{osmol/L}$（\qty{100}{g}
$= \qty{0.55}{mol}$ 游离葡萄糖）。
\textbf{19.} 几乎是\omterm{def:b1:eukaryotic-cell:organelle}{细胞溶胶}渗透浓度的两倍：水会从血液中涌入，
\omterm{def:b1:cell-unit-of-life:cell}{细胞}将膨胀至接近三倍的
体积并破裂。
\textbf{20.} $\num{7e18}/\num{6e23} = \qty{1.2e-5}{mol}$ 溶于 \qty{1}{L} 中：
\qty{12}{\micro osmol/L}，少了四万倍。
\textbf{21.} $\Psi_s = -2.58\times 0.55 = \qty{-1.4}{MPa}$：\omterm{def:b1:eukaryotic-cell:plantcell}{细胞壁}
需要承受 \qty{1.4}{MPa}，即十四个大气压。
\textbf{22.} 肌\omterm{def:b1:carbohydrates:polysaccharide}{糖原}是肌肉自身收缩所用的燃料，
以葡萄糖-6-磷酸的形式直接动员进入糖酵解；肌肉
缺乏能将葡萄糖从磷酸基团上释放出来的酶，因此该储备物
无法离开肌纤维。
\textbf{23.} 马铃薯在发芽的数周内动员其\omterm{def:b1:carbohydrates:polysaccharide}{淀粉}；肝脏则必须
在数分钟内作出反应。大而致密、表面末端较少的颗粒
适合缓慢而稳定的释放；小而高度分支、带有数千个末端的粒子
则适合快速释放。
\textbf{24.} $0.62/(4\times 3600) = \qty{4.3e-5}{mol/s} = \num{2.6e19}$
残基每秒；\omterm{def:b1:water-small-molecules:atp}{ATP} 合计 \qty{1.24}{mol}，每秒
$\num{5.2e19}$ 个。
\textbf{25.} 释放：每秒 $\num{9.3e18}$ 个葡萄糖分子
（\qty{10}{g/h}）；由 \num{1.4e22} 个非还原末端实现，
每个粒子 \num{2048} 个；这种多聚体使\omterm{def:b1:cell-unit-of-life:cell}{细胞}免于
承受 \qty{0.55}{osmol/L} 的渗透压，该值几乎是\omterm{def:b1:cell-unit-of-life:cell}{细胞}自身渗透浓度的两倍。
\end{solution}
